\chapter{Volumes}

%=============================================================================
\section{Volumes by Slicing}
\label{sec:10.1}
%=============================================================================

\textit{We defined integrals in order to compute areas, but there are many other applications of integrals.  In this section we revisit the definition of the Riemann integral, reinterpret the notation with a different geometrical flavour, and use the same idea to set up volumes as integrals.}

\subsection*{Reinterpreting the integral}

Recall the idea behind the Riemann integral.  Given a region bounded by the graph of a function $f$, the $x$-axis, and the vertical lines $x = a$ and $x = b$, we break $[a,b]$ into smaller pieces using a partition.  This cuts the region into thin vertical slices.  Each slice is approximated by a rectangle of width $\Delta x$ and height $f(x)$, so its area is approximately $f(x)\cdot\Delta x$.  Summing over all slices gives an approximation to the total area, and taking the limit as $\Delta x \to 0$ yields the definite integral.

\textit{What follows is not a fully rigorous argument, but it provides powerful geometric intuition and helps us reuse the same ideas for many applications of the integral beyond areas.}

Since the integral is a limit of sums, we abuse notation and think of it as an actual sum of infinitely many, infinitesimally small pieces.  Specifically, we write $\dx$ and think of it as a $\Delta x$ that is \textbf{infinitesimally small}.  The integral sign then represents the ``sum'' over infinitely many such pieces.

Consider the same region as before.  Instead of finitely many slices, imagine breaking it into infinitely many slices, each of infinitesimal width $\dx$.  Focus on a single slice.  Its width is $\dx$, and since the slice is infinitesimally thin, its area is also infinitesimal; call it $dA$.

\begin{remark}
When we approximate such a slice by a rectangle, the error is on the order of $(\dx)^{2}$: it is infinitesimally small \textit{squared}.  Compared with the slice itself, this error is negligible, and we may ignore it in the limit.
\end{remark}

Therefore the area of the slice is
\[
  dA = f(x)\dx,
\]
and the total area is the ``sum'' of all such infinitesimal pieces:
\[
  A = \int_{a}^{b} f(x)\dx.
\]

\subsection*{From areas to volumes}

\textit{The advantage of thinking of integrals this way is that it helps with geometric intuition, and it helps us reuse the same ideas for many other applications --- for example, computing volumes.}

Suppose we have a solid (picture a carrot) whose volume we wish to compute.  If we approximate it crudely by a single cylinder, the result is poor.  But just as we sliced a region into thin rectangles for areas, we can \textbf{slice the solid} into thin pieces.  Each slice is approximately a disc, and we know how to compute the volume of a disc.  Taking the limit as the thickness tends to zero, the approximation becomes exact.

\begin{definition}[Volume by slicing --- disc method]
\label{def:10-1}
Let a solid be generated by rotating the region between the graph of a function $y = f(x)$ and the $x$-axis, from $x = a$ to $x = b$, around the $x$-axis.  Cut the solid into infinitesimally thin slices perpendicular to the $x$-axis.  Each slice is a disc of radius $R = f(x)$ and thickness $\dx$.  Its infinitesimal volume is
\[
  dV = \pi R^{2}\dx = \pi\bigl[f(x)\bigr]^{2}\dx,
\]
and the total volume is
\[
  V = \int_{a}^{b} \pi\bigl[f(x)\bigr]^{2}\dx.
\]
\end{definition}

% ── BEGIN VISUAL ──
\begin{figure}[ht]
  \centering
  \begin{tikzpicture}
  \begin{axis}[
    width=11cm,
    height=6.8cm,
    axis lines=middle,
    xmin=-0.2, xmax=4.4,
    ymin=-0.4, ymax=2.6,
    xlabel={$x$}, ylabel={$y$},
    ticks=none,
    clip=true
  ]
    % Region under y = 2 - x/2 on [0,4], rotated about x-axis
    \addplot[thick, color=thmcolor, domain=0:4, samples=2] {2 - 0.5*x};
    \addplot[domain=0:4, draw=none, fill=excolor, fill opacity=0.12] {2 - 0.5*x} \closedcycle;

    % Representative slice at x=2
    \def\xS{2}
    \pgfmathsetmacro{\rS}{2 - 0.5*\xS}
    \draw[thick, color=defcolor] (axis cs:\xS,0) -- (axis cs:\xS,\rS);
    \draw[thick, color=defcolor] (axis cs:\xS-0.12,0) -- (axis cs:\xS+0.12,0);
    \node[color=defcolor, anchor=west] at (axis cs:2.15,\rS/2) {$r=f(x)$};

    \node[anchor=west, font=\small] at (axis cs:0.3,2.35) {rotate about $x$-axis $\Rightarrow$ disk};
  \end{axis}
\end{tikzpicture}

  \caption{Disk method intuition: a slice rotates into a disk.}
  \label{fig:disk-method-slice}
\end{figure}
% ── END VISUAL ──

\begin{remark}
Just as with areas, the key step is identifying the infinitesimal piece (a thin disc), writing its volume in terms of the variable of integration, and then summing (integrating) over all such pieces.
\end{remark}

\subsection*{A concrete example}

\begin{example}
\label{ex:10-1}
Let $R$ be the region in the first quadrant between the graph of $y = x^{2}$ and the $x$-axis, from $x = 0$ to $x = 1$.  Rotate $R$ around the $x$-axis.  Compute the volume of the resulting solid.
\end{example}

\begin{remark}[Pause and Try]
Before reading the solution, draw a picture and try to understand what the solid looks like.  It resembles a cone that has been bent slightly.
\end{remark}

\textbf{Solution.}  Slice the solid into infinitesimally thin pieces perpendicular to the $x$-axis.  A single slice is obtained by taking a thin vertical strip of width $\dx$ from the region $R$ and rotating it around the $x$-axis.  This produces a disc.

The disc has thickness $\dx$ and radius equal to the height of the strip, which is $y = x^{2}$.  Therefore the infinitesimal volume is
\[
  dV = \pi\bigl(x^{2}\bigr)^{2}\dx = \pi x^{4}\dx.
\]
Integrating from $x = 0$ to $x = 1$:
\[
  V = \int_{0}^{1} \pi x^{4}\dx = \pi\cdot\frac{x^{5}}{5}\bigg|_{0}^{1} = \frac{\pi}{5}.
\]

\begin{remark}
The strategy is always the same: slice the solid into infinitesimally thin pieces, write the volume of each piece, and integrate.  The calculations may be more involved for other solids, but the underlying idea is identical --- slice the carrot infinitesimally thin.
\end{remark}

%=============================================================================
\section{Cylindrical Shells}
\label{sec:10.2}
%=============================================================================

\textit{In the previous section, we computed volumes by slicing a solid perpendicular to the axis of revolution, producing discs (or washers).  In this section, we work with a different example that requires a different approach.  The key idea is still slicing the solid into infinitesimally thin pieces; we simply slice in a different direction.}

\subsection*{When the disc method fails}

\begin{example}
\label{ex:10-2}
Let $R$ be the region in the first quadrant between the graph of $y = 2x^{4} - x^{5}$ and the $x$-axis.  Rotate $R$ around the $y$-axis.  Compute the volume of the resulting solid.
\end{example}

\begin{remark}[Pause and Try]
Before reading on, draw a picture and try to visualise the solid.  It looks like a small volcano --- a mountain with a hole in the centre.  The curve $y = 2x^{4} - x^{5}$ forms the profile, and rotating it around the $y$-axis produces the volcano shape.
\end{remark}

\textit{Let us first attempt the approach from \cref{sec:10.1}: slice perpendicular to the axis of revolution (the $y$-axis).  This means slicing horizontally.}

A thin horizontal strip of height $dy$ in the region $R$ corresponds, after rotation, to a \emph{washer}: a disc with a hole.  If the inner radius is $r$ and the outer radius is $R$, the infinitesimal volume would be
\[
  dV = \pi\bigl(R^{2} - r^{2}\bigr)\,dy.
\]
To finish the setup, we would integrate with respect to $y$ from $0$ to the maximum value of the function, and we would need to express $R$ and $r$ in terms of $y$.  That means solving
\[
  y = 2x^{4} - x^{5}
\]
for $x$ in terms of $y$.

\begin{remark}[Warning]
This equation cannot be solved for $x$ by elementary methods.  The plan was a good idea and worth trying, but ultimately we cannot complete the problem this way.  We need a different approach.
\end{remark}

\subsection*{The cylindrical shell method}

\textit{Instead of slicing horizontally (perpendicular to the $y$-axis), we slice the region \textbf{vertically} (parallel to the $y$-axis) and integrate with respect to $x$.}

Take a thin vertical strip of width $\dx$ in the region $R$.  When this strip is rotated around the $y$-axis, it sweeps out a \emph{cylindrical shell}: a thin hollow cylinder.

\begin{definition}[Volume by cylindrical shells]
\label{def:10-2}
When a thin vertical strip at position $x$, of height $y = f(x)$ and width $\dx$, is rotated about the $y$-axis, it produces a cylindrical shell of
\begin{itemize}
  \item \textbf{radius} $x$ (the distance from the strip to the axis),
  \item \textbf{height} $y = f(x)$,
  \item \textbf{thickness} $\dx$.
\end{itemize}
Because the thickness is infinitesimally small, we may ``unroll'' the shell into a thin rectangular brick of dimensions $2\pi x \times f(x) \times \dx$.  The infinitesimal volume is therefore
\[
  dV = 2\pi x\cdot f(x)\dx,
\]
and the total volume is
\[
  V = \int_{a}^{b} 2\pi x\cdot f(x)\dx.
\]
\end{definition}

% ── BEGIN VISUAL ──
\begin{figure}[ht]
  \centering
  \begin{tikzpicture}
  \begin{axis}[
    width=11cm,
    height=6.8cm,
    axis lines=middle,
    xmin=-0.2, xmax=4.4,
    ymin=-0.4, ymax=3.0,
    xlabel={$x$}, ylabel={$y$},
    ticks=none,
    clip=true
  ]
    % Region under y = 2 - x/2 on [0,4], rotated about y-axis (shells)
    \addplot[thick, color=thmcolor, domain=0:4, samples=2] {2 - 0.5*x};
    \addplot[domain=0:4, draw=none, fill=excolor, fill opacity=0.12] {2 - 0.5*x} \closedcycle;

    % Representative vertical strip at x=2
    \def\xS{2}
    \pgfmathsetmacro{\hS}{2 - 0.5*\xS}
    \path[fill=defcolor, fill opacity=0.18] (axis cs:\xS-0.12,0) rectangle (axis cs:\xS+0.12,\hS);
    \draw[thick, color=defcolor] (axis cs:\xS-0.12,0) rectangle (axis cs:\xS+0.12,\hS);

    % Axis of rotation (y-axis)
    \draw[very thick, color=black] (axis cs:0,-0.4) -- (axis cs:0,3.0);

    \node[anchor=west, font=\small] at (axis cs:0.3,2.7) {rotate about $y$-axis $\Rightarrow$ shell};
    \node[color=defcolor, anchor=west] at (axis cs:2.15,1.0) {height $=f(x)$};
  \end{axis}
\end{tikzpicture}

  \caption{Shell method intuition: a strip rotates into a shell.}
  \label{fig:cylindrical-shells}
\end{figure}
% ── END VISUAL ──

\begin{remark}
Strictly speaking, unrolling the shell into a brick is an approximation.  But because the thickness $\dx$ is infinitesimally small, the error from flattening the curvature is negligible --- it is of order $(\dx)^{2}$.  If the thickness were not infinitesimal, the error could be significant.
\end{remark}

\subsection*{Completing the example}

We return to \cref{ex:10-2}.  Slice the region $R$ vertically.  A thin strip at position $x$ has width $\dx$ and height $y = 2x^{4} - x^{5}$.  Rotating this strip around the $y$-axis produces a cylindrical shell of radius $x$, height $2x^{4} - x^{5}$, and thickness $\dx$.  Its infinitesimal volume is
\[
  dV = 2\pi x\bigl(2x^{4} - x^{5}\bigr)\dx.
\]
Since $y = 2x^{4} - x^{5} = x^{4}(2 - x)$, the function is zero at $x = 0$ and $x = 2$, so we integrate from $0$ to $2$:
\[
  V = \int_{0}^{2} 2\pi x\bigl(2x^{4} - x^{5}\bigr)\dx = 2\pi\int_{0}^{2}\bigl(2x^{5} - x^{6}\bigr)\dx.
\]
This integral can now be evaluated using antiderivatives.

\subsection*{Comparing the two methods}

\textit{Both methods are fundamentally the same idea: cut the solid into infinitesimally small pieces and integrate.  We simply happen to cut in different directions.}

\begin{definition}[Washer method vs.\ cylindrical shell method]
\label{def:10-3}
\leavevmode
\begin{itemize}
  \item The \emph{washer method} (or disc method): we cut the region \textbf{perpendicular} to the axis of revolution.  Each piece is a disc or washer.
  \item The \emph{cylindrical shell method}: we cut the region \textbf{parallel} to the axis of revolution.  Each piece is a cylindrical shell.
\end{itemize}
\end{definition}

% ── BEGIN VISUAL ──
\begin{figure}[ht]
  \centering
  \begin{tikzpicture}
  \begin{axis}[
    width=11cm,
    height=6.8cm,
    axis lines=middle,
    xmin=-0.2, xmax=4.4,
    ymin=-0.4, ymax=3.0,
    xlabel={$x$}, ylabel={$y$},
    ticks=none,
    clip=true
  ]
    % Region between y=2 and y=2-x/2 on [0,4], rotated about x-axis
    \addplot[thick, color=thmcolor, domain=0:4, samples=2] {2};
    \addplot[thick, color=thmcolor, domain=0:4, samples=2] {2 - 0.5*x};

    % Fill the region
    \addplot[draw=none, fill=excolor, fill opacity=0.12, domain=0:4] {2} \closedcycle;
    \addplot[draw=none, fill=white, domain=0:4] {2 - 0.5*x} \closedcycle;

    % Representative slice at x=2
    \def\xS{2}
    \pgfmathsetmacro{\rOut}{2}
    \pgfmathsetmacro{\rIn}{2 - 0.5*\xS}
    \draw[thick, color=defcolor] (axis cs:\xS,0) -- (axis cs:\xS,\rOut);
    \draw[thick, color=defcolor] (axis cs:\xS,0) -- (axis cs:\xS,\rIn);
    \node[color=defcolor, anchor=west] at (axis cs:2.15,1.7) {$R$};
    \node[color=defcolor, anchor=west] at (axis cs:2.15,0.7) {$r$};

    \node[anchor=west, font=\small] at (axis cs:0.3,2.7) {rotate $\Rightarrow$ washer ($\pi(R^2-r^2)$)};
  \end{axis}
\end{tikzpicture}

  \caption{Washer method intuition: outer minus inner disk.}
  \label{fig:washer-method-slice}
\end{figure}
% ── END VISUAL ──

\begin{remark}
The names are less important than the underlying principle.  In any given problem, one method may lead to easier integrals than the other.  When one approach produces an equation that cannot be solved, try slicing in the other direction.
\end{remark}
